% ==========================================================================
% GENERATED by scripts/split_paper.py. Do not edit this file.
%
% Source: imported_papers/Neural_Computation__CNA/NECO-D-24-00304-Source.tex
% Part 4: Experimental Results
%
% Edits belong upstream, in the published paper. To change what the
% chapter includes, change imports/cna.toml and re-run
% the splitter.
% ==========================================================================
\section{Experimental Results}\label{sec:exp_results}

\subsection{Experimental Setup}\label{sec:exp_setting}
Our experiments are designed to rigorously probe the \emph{core representational capability} of the CNA: its ability to form net fragments during training and recombine them into larger nets during inference. To isolate this property, we deliberately avoid task-specific or dataset-driven complexity and instead focus on the model’s behavior in its most fundamental setting.

We therefore train CNA on binary images of straight lines at varying orientations. Straight lines serve as the most elementary visual primitives: they are simple enough to avoid confounds yet naturally compose into richer structures such as line drawings of objects. This controlled regime directly tests whether CNA can learn fragmentary representations and later reassemble them into unseen global patterns. By contrast, benchmarks such as MNIST \cite{lecun1998gradient}, CIFAR \cite{krizhevsky2009cifar}, or ImageNet \cite{Deng_2009} involve dense, texture-heavy images where the link between local features and global structure is opaque. Evaluating CNA on such data at this stage would obscure its core representational mechanism, and, as we discuss in Section~\ref{sec:discussion}, would require introducing an additional object-representing area and more advanced filters, which we defer to future work.

Evaluation targets two key properties:
\begin{itemize}
    \item \textbf{Compositionality:} Can learned line fragments recombine into novel structures, such as object line drawings unseen during training?
    \item \textbf{Robustness:} Can nets be retrieved reliably under noise or occlusion?
\end{itemize}
To this end, we present the model \emph{after training and without updating any parameters} with (i) line drawings of objects composed of straight-line segments, (ii) inputs corrupted with Gaussian noise, and (iii) partially occluded patterns.

For comparison, we also train a conventional autoencoder on the same line-stimuli dataset to highlight representational differences. Further experimental details are provided in the appendix: We describe the training and testing datasets in more detail in Appendix \ref{sec:dataset}, the feature extractor mechanism to obtain $\boldsymbol{y}^{(S1)}$ in Appendix \ref{sec:features}, the training details of our CNA model in Appendix \ref{sec:parameters}, and the training details of the autoencoder in Appendix \ref{sec:autoencoder_params}.


\subsection{Interpreting Figures: General Considerations}

\begin{figure}[ht!]
  \begin{center}
    \includegraphics[width=\textwidth]{NECO-D-24-00304-Figure.3}
  \end{center}
  \caption{%
  \textbf{The response of our model to a diagonal line input} is displayed as a net of activated neurons. The left of the image shows the neuronal activation $\boldsymbol{y}^{(S2)}$ when observing the line, and the right visualizes the net formed by these neurons (activations $\boldsymbol{y}^{(S2)}$ with their connections $\boldsymbol{W}^{(L)}$). For convenience, the neurons are rearranged in a circle, and the net fragment (connections) of two randomly selected neurons, $n_i$ and $n_j$, are visualized in colors (purple and red).}\label{fig:result_interpretation}
\end{figure}

A defining feature of CNA is that it represents patterns not through individual neuron activations, but through \emph{nets} -- cohesive, interconnected clusters of mutually supportive neurons. This network-based organization means that an individual neuron's activation is meaningful only in the context of its embedding within a larger net. Understanding CNAs' outputs, therefore, requires a shift in perspective: from isolated neuron activity to the connectivity and structure of the nets they form.

Figure \ref{fig:result_interpretation} illustrates this principle. After presenting a diagonal line input, the left panel shows the resulting neuronal activity $\boldsymbol{y}^{(S2)}$, while the right panel visualizes the underlying net structure derived from recurrent connections $\boldsymbol{W}^{(L)}$. For clarity, neurons are arranged in a circular layout, and the net fragments $\mathcal{F}(n_i)$ and $\mathcal{F}(n_j)$ of two randomly selected neurons $n_i$ and $n_j$ are highlighted in color. 

This visualization demonstrates two important points. First, individual neurons (e.g., $n_i$ and $n_j$) remain active only if embedded within supportive nets: without sufficient recurrent support, they are suppressed by attenuation mechanisms and would appear inactive (black) in the left panel. Second, this shows that observed activity patterns should always be interpreted as the visible footprint of underlying nets, even when those nets are not explicitly drawn.

Consequently, all subsequent visualizations of neural activity must be understood in this net-based context: the squares represent neurons that survive attenuation precisely because they are supported by their respective nets.

All our visualizations of neural activity are color-coded according to the competing channel groups $\mathcal{C}(n_i)$. Each color corresponds to a feature type (e.g., horizontal, vertical, or diagonal lines; see Appendix~\ref{sec:features}). However, these colors do not distinguish between different competitive channels $k \in {1, \dots, \kappa }$ within the same feature group, meaning that two neurons with the same color may belong to different competitive subnetworks with distinct connectivity patterns.


\subsection{Compositionality}\label{sec:results_compositionality}
%
\begin{figure}[ht!]
  \begin{center}
    \includegraphics[width=\textwidth]{NECO-D-24-00304-Figure.4}
  \end{center}
  \caption{%
  \textbf{Training and evaluation data.} The top of the figure illustrates the straight lines used during the training phase. Below, various samples from the evaluation datasets are shown, including kinked lines (rows $1$ and $2$), digits and characters (rows $3$ and $4$), and line drawings (rows $5$--$8$). For each row, two input examples, denoted as $\boldsymbol{x}_1$ and $\boldsymbol{x}_2$, are displayed alongside their respective feature activations, $\boldsymbol{y}^{(S1)}$, and neuronal activations, $\boldsymbol{y}^{(S2)}$, shown both without and with addition of Gaussian noise.
  }\label{fig:compositionality}
\end{figure}
%
The model weights are trained using straight lines (see training data examples at the top of Figure \ref{fig:compositionality}). During inference -- without additional learning -- the learned net fragments are recursively composed to form nets representing significantly more complex structures (see validation data in rows $1$--$8$ of Figure \ref{fig:compositionality}).
Notably, although such structures are absent from the training set, the CNA successfully generalizes to represent kinked lines (rows $1$ and $2$), digits and characters (rows $3$ and $4$), and line drawings (rows $5$--$8$).

These results demonstrate that while individual net fragments  $\mathcal{F}_t$ must be learned (in our case we use Hebbian learning to train the connections within $\mathcal{F}_t$), the recursive composition described in Eq. \ref{eq:recursive_net_fragments} to form $\mathcal{N}^*_t$ allows the integration of previously uncombined fragments. Consequently, the model exhibits combinatorial generalization through the structural recomposition of learned primitives.

\subsection{Filtering Gaussian Noise}\label{sec:results_noise_filtering}


The motivation for why we construct net fragments rather than just using the feature activation $\boldsymbol{y}^{(S1)}$ is demonstrated by introducing additional clutter:
For each sample, we generate an input representation both with and without Gaussian noise. This noise corresponds to invalid patterns not observed during training (i.e., are not captured in net fragments) that consequently lack sufficient recurrent support, leading to the deactivation of the corresponding neurons.
Thus, the CNA model only allows valid sub-patterns to persist and be composed into novel structures, while invalid patterns, such as noise, are suppressed.

\begin{figure}[ht!]
  \begin{center}
    \includegraphics[width=\textwidth]{NECO-D-24-00304-Figure.5}
  \end{center}
  \caption{\textbf{Noise filtering efficacy}: Noise-corrupted feature activations $\boldsymbol{y}^{(S1)}$ alongside the corresponding output $\boldsymbol{y}^{(S2)}$, in which noise is reduced due to a lack of support within nets. %The plots are created using a bias $b^{(S2)} = 0.5$ and an attenuation coefficient $\gamma=1.2+0.2t$.
  }\label{fig:noise_filtering}
\end{figure}
To further investigate noise filtering, we introduce Gaussian noise of varying degrees (from $0\%$ to $59\%$ noise at each spatial location - resulting from adding up to $20\%$ noise per channel across four independent channels; see Appendix \ref{sec:dataset}) to the afferent input in stage $S2$, thereby simulating an image structure that has already survived the filtering of feature kernels in $S1$ (thus increasing the difficulty, as the fragments must handle all the noise rather than relying on the feature extraction stage).

The results are visualized in Figure \ref{fig:noise_filtering}: While introducing up to $47.8\%$ noise, remarkably consistent outputs $\boldsymbol{y}^{(S2)}$ are generated. As additional noise is introduced, more undesired neurons receive support and persist in their activity. However, the overarching pattern remains distinguishable even when subjected to noise levels of up to $59\%$.
In Appendix \ref{sec:quant_noise_filtering}, we quantitatively confirm the efficacy of noise filtering by measuring precision, recall, and noise filtration rate for various parameter settings, showing that lines remain well distinguishable for up to $59\%$ added noise.

\begin{figure}[ht!]
  \begin{center}
    \includegraphics[width=\textwidth]{NECO-D-24-00304-Figure.6}
  \end{center}
  \caption{\textbf{The filtering of (simulated) noise over time}: Already at time step $t=0$, most of the noise is suppressed. However, when attenuation increases over time, even more noise is filtered out. The plot is created after adding $10\%$ noise to each feature channel (using a bias $b^{(S2)} = 0.7$).% and an attenuation coefficient $\gamma=1.2+0.2t$).
  }\label{fig:noise_time}
\end{figure}

The reason that the filtering works is that the additionally introduced clutter doesn't activate overlapping net fragments and cannot persist over time.
Figure \ref{fig:noise_time} depicts this systematic filtering across timesteps.
The figure illustrates a substantial noise reduction already at time step $t=0$ due to initial attenuation. In each subsequent time step, additional clutter is removed. This reduction is due to increased attenuation and falling support as a consequence of the drop-out of supporting neurons.


\subsection{Reconstructing Figures}\label{sec:partial_occlusion}

\begin{figure}[ht!]
  \begin{center}
    \includegraphics[width=0.95\textwidth]{NECO-D-24-00304-Figure.7}
  \end{center}
  \caption{%
  \textbf{Recovery from occlusion.} 
  Part A: %
  Input images of diagonal lines with different numbers of missing pixels, corresponding feature activations (middle column), and final output neuron activations. Activation bias $b^{(S2)}=0.5$, attenuation coefficient $\gamma=0.6+0.2t$. %
  %
  Part B: %
  A diagonal line with $6$ missing pixels in the center (shown on the left), the corresponding output activity (in the middle), and the pre-synaptic activity $\boldsymbol{a}^{(S2)}$ displaying the ``activation probabilities'' (on the right) after the first time step.}
  \label{fig:li_examples}
\end{figure}
If a net fragment $\mathcal{F}_t(n_i)$ contains a substantial number of (active) neurons, an inactive neuron $n_i$ can receive significant support, leading it to switch on and encouraging figure reconstruction ($\mathcal{F}_t(n_i) \rightarrow \mathcal{F}^*_t(n_i)$).
Figure \ref{fig:li_examples}A shows how inactive neurons activate and thus demonstrates that nets can deal with occluded patterns.
This reconstruction is reliable for up to $3$ missing pixels (see Appendix \ref{sec:quant_noise_filtering}).
In some cases, as in this figure, the line is fully reconstructed for up to $6$ removed pixels and partly reconstructed for $7$ pixels.

\paragraph{Pre-Synaptic Activity Represents Expected Feature Activation.}
In the presence of missing pattern elements, the pre-synaptic activity map $\boldsymbol{a}^{(S2)}$ is higher at spatial locations where line features are observed and lower where they are not observed but expected (as shown in Figure \ref{fig:li_examples}B).
Depending on the context, it may be desirable either to reconstruct the missing line segment -- effectively bridging the gap between discontinuous features -- or to preserve the separation between distinct line elements.
This behavior can be modulated by adjusting the hyperparameter $b^{(S2)}$ in Eq. \ref{eq:s2_final_output}, which can govern the balance between feature completion and preservation of discontinuity, thereby allowing the model to selectively favor line reconstruction or separation. 
In future work, incorporating additional areas that maintain object-associated representations and act as top-down hypotheses could enable this selectivity to be learned rather than manually tuned.

\subsection{Comparison to Autoencoders}\label{sec:results_early_commitment}
\begin{figure}[ht!]
  \begin{center}
    \includegraphics[width=\textwidth]{NECO-D-24-00304-Figure.8}
  \end{center}
  \caption{\textbf{Comparison of the outputs generated by the proposed CNA model and an autoencoder}. Each row presents two samples, $\boldsymbol{x}_1$ and $\boldsymbol{x}_2$, alongside their respective feature activations and resulting representations (once for the CNA model and once for the autoencoder). Rows $1$--$3$ illustrate how both models encode patterns that were not present in the training dataset. Rows $4$--$5$ let one compare the models' performance in denoising tasks, while row $6$ displays their ability to reconstruct lines.}
  \label{fig:results_early_commitment}
\end{figure}

Figure \ref{fig:results_early_commitment}  compares performance between the proposed CNA model and an autoencoder (for model and training details, refer to Appendix \ref{sec:autoencoder_params}). While the autoencoder successfully learns to accurately represent line structures observed in the training data -- achieving a mean square error (MSE) of less than $1 \times 10^{-3}$ between the input and the reconstruction by the decoder -- it struggles with generalization to data not observed during training.
For instance, rows $1$--$3$ demonstrate that the CNA model is capable of representing previously unseen structures, while, in contrast, the autoencoder fails to capture such out-of-distribution patterns.

Rows $4$--$5$ highlight the autoencoder's ability to represent line structures with up to 11.5\% noise introduced. However, as the noise level increases, its performance decreases significantly. The CNA model, in contrast, remains robust under much higher noise levels, tolerating up to $59\%$ noise while still producing high-quality representations. Metrics on noise reduction, precision, and recall, presented in Appendix \ref{sec:quant_noise_filtering}, further substantiate the CNA model's superior performance.

In row $6$, the models' capabilities for reconstructing missing line segments are compared. Both models can reconstruct the removed pixels, but the autoencoder introduces additional noise in regions other than the line center, where the pixels were removed.
This observation is confirmed by the metrics in Appendix \ref{sec:quant_noise_filtering}: The autoencoder is slightly superior in reconstructing the removed pixels, but this comes at the cost of increased noise elsewhere, leading to lower overall precision and recall for figure reconstruction.
